\documentclass[11pt,reqno]{amsart}
\usepackage[T1]{fontenc}
\usepackage[utf8]{inputenc}
\usepackage{lmodern}
\usepackage{amsmath,amssymb,amsfonts,amsthm,mathtools,mathrsfs}
\usepackage{enumitem}
\usepackage{microtype}
\usepackage[colorlinks=true,linkcolor=blue,citecolor=blue,urlcolor=blue]{hyperref}
\usepackage[nameinlink,capitalize]{cleveref}
\allowdisplaybreaks
\setlength{\emergencystretch}{2em}

\newtheorem{theorem}{Theorem}[section]
\newtheorem{proposition}[theorem]{Proposition}
\newtheorem{lemma}[theorem]{Lemma}
\newtheorem{corollary}[theorem]{Corollary}
\theoremstyle{definition}
\newtheorem{definition}[theorem]{Definition}
\newtheorem{assumption}[theorem]{Assumption}
\newtheorem{example}[theorem]{Example}
\theoremstyle{remark}
\newtheorem{remark}[theorem]{Remark}

\newcommand{\A}{\mathcal A}
\newcommand{\B}{\mathcal B}
\newcommand{\C}{\mathbb C}
\newcommand{\E}{\mathbb E}
\newcommand{\N}{\mathbb N}
\newcommand{\R}{\mathbb R}
\newcommand{\T}{\mathbb T}
\newcommand{\X}{\mathcal X}
\newcommand{\F}{\mathcal F}
\newcommand{\Piop}{\Pi}
\newcommand{\Id}{\mathrm{Id}}
\newcommand{\Var}{\operatorname{Var}}
\newcommand{\Tot}{\operatorname{Tot}}
\newcommand{\supp}{\operatorname{supp}}
\newcommand{\norm}[1]{\left\lVert #1\right\rVert}
\newcommand{\abs}[1]{\left\lvert #1\right\rvert}
\newcommand{\pair}[2]{\left\langle #1,#2\right\rangle}

\title[Bilateral response and symbolic desingularization]{Bilateral Response, Symbolic Desingularization, and All-Order Spectral Jets for Nonconjugate Moving Collision Dynamics}
\author{Qian Qi}
\address{Peking University, Beijing 100871, China}
\email{qiqian@pku.edu.cn}
\subjclass[2020]{Primary 37C30, 37D50; Secondary 37A25, 47A55, 60G42}
\keywords{linear response, moving singularities, graph currents, open dispersing billiards, Ruelle operators, higher response}
\date{August 29, 2026}

\begin{document}

\begin{abstract}
We prove positive all-order response theorems for two genuinely nonconjugate
moving collision systems.  The first is an explicit symplectic Markov collision
map with three moving seams.  Its unstable quotient is a four-branch full-cover
map.  On a fixed Sobolev--variation ladder the centered transfer operator
contracts uniformly, every parameter derivative is explicit, arbitrary
independently prescribed correlation susceptibilities are differentiable to
every finite order, and the same nonautonomous map produces exact predictable
innovations.  No assembly reset or exact-coboundary ansatz is used.

The second is a family of analytic open dispersing billiards satisfying the
no-eclipse condition.  A parameter-dependent graph transform gives an analytic
coding map on one fixed subshift of finite type.  Pulling roof functions,
unstable Jacobians, and physical observables to that fixed symbolic space
removes the moving singularity before differentiation.  The resulting Ruelle
operators form an analytic family on one Hölder space.  Pressure, equilibrium
states, arbitrary Hölder correlation susceptibilities, reduced resolvents, and
all finite spectral jets are therefore analytic.  Moving one obstacle changes
an isolated two-periodic flight length, so the family is not time-preservingly
conjugate.

The common mechanism is a sharp bilateral product-tail theorem.  Source-side
and target-side crossed estimates interpolate to an absolutely summable
product envelope exactly when a strict rate determinant is positive.  We also
prove finite-difference convergence in the complete $\ell^1(\N^2)$ response
topology and a joint-Cauchy theorem for totalizing higher sources on a graded
Banach scale.
\end{abstract}

\maketitle
\tableofcontents

\section{Introduction}

Differentiating a smooth operator family on one Banach space is governed by
classical perturbation theory \cite{Kato1995}.  A moving collision boundary is
harder: before physical assembly the derivative is a signed current carried by
a moving face, and after propagation it remembers a waiting time on each side
of the insertion.  One-sided decay does not sum the resulting two-time array.

This paper gives two positive ways to close the problem.  The first retains the
physical seam currents and proves a bilateral product estimate directly.  The
second creates a fixed symbolic coordinate system for an open dispersing
billiard; in that coordinate system the moving obstacle geometry appears as an
analytic Hölder potential rather than as a moving discontinuity.  We call this
second tool \emph{symbolic desingularization}.

The explicit Markov collision system also closes the probabilistic interface
needed downstream.  Even if its parameter at time $k$ is selected from the
previous collision history, the present unstable coordinate is conditionally
Lebesgue, the next branch has the declared parameter-dependent probability,
and the image coordinate is uniform and independent of the enlarged past.
Thus the martingale arrays in the homogenization paper arise from the same
deterministic nonautonomous collision map.

Our functional-analytic estimates are compatible with strong--weak spectral
perturbation \cite{KellerLiverani1999}, response for piecewise expanding maps
\cite{BaladiSmania2008}, and anisotropic methods for hyperbolic systems with
singularities \cite{DemersLiverani2008,DemersZhang2011}.  The open-billiard
branch uses the standard symbolic coding of no-eclipse dispersing billiards and
Ruelle--Perron--Frobenius theory \cite{Bowen1975,ParryPollicott1990,
ChernovMarkarian2006}.  The new point is that the coding is differentiated as a
map into a fixed Hölder space, yielding an actual moving-scatterer response
theorem rather than a fixed-table stability statement.

\section{Occurrencewise bilateral response}\label{sec:bilateral}

Let $a$ range over a compact interval $\A$.  Let $L_a$ be a transfer operator
with a simple invariant projection $\Pi_a$ and centered part
\[
 Q_a=L_a-\Pi_a.
\]
All fibres are transported to a common measurable reference space.

\begin{definition}[Same-occurrence registry]\label{def:registry}
An occurrence label $\omega$ contains the physical face, the incident side,
the owner branch, the oriented branch word, the parameter chamber, and the
refinement chart.  Labels are identified only on a common refinement
preserving all these data.  Reversing orientation changes the sign but not the
UID.
\end{definition}

Write the physical first derivative as
\begin{equation}\label{eq:source}
 \partial_aL_a=K_a^{\rm reg}+\sum_{\omega\in\Omega_a}J_{a,\omega}.
\end{equation}
For centered $g$ and an admissible test $f$, set
\begin{equation}\label{eq:atom}
 A_{m,n}^{\omega}(a;g,f)
 =\pair{Q_a^nJ_{a,\omega}Q_a^m g}{f},
 \qquad m,n\ge0.
\end{equation}
No cancellation between distinct UIDs is used before a common-refinement
identity has been proved.

\begin{assumption}[Crossed envelopes]\label{ass:crossed}
There are $0<\alpha,\beta<1$, $\Gamma_+,\Gamma_-\ge1$, and constants
$C_\omega^\pm\ge0$ such that
\begin{align}
 \abs{A_{m,n}^{\omega}}
 &\le C_\omega^-\alpha^m\Gamma_+^n
       \norm g_{X_-}\norm f_Y,\label{eq:source-side}\\
 \abs{A_{m,n}^{\omega}}
 &\le C_\omega^+\Gamma_-^m\beta^n
       \norm g_{X_-}\norm f_Y.\label{eq:target-side}
\end{align}
\end{assumption}

Put
\[
 A_0=\log(1/\alpha),\qquad B_0=\log(1/\beta),\qquad
 C_0=\log\Gamma_+,\qquad D_0=\log\Gamma_-.
\]

\begin{theorem}[Sharp crossed-rate product tail]\label{thm:crossed}
Assume \cref{ass:crossed}.  There is an interpolation parameter for which both
time factors are strictly smaller than one if and only if
\begin{equation}\label{eq:determinant}
 A_0B_0>C_0D_0.
\end{equation}
Under this condition define
\[
 \lambda_*={B_0+D_0\over A_0+B_0+C_0+D_0},
 \qquad
 \log\rho_*={C_0D_0-A_0B_0\over A_0+B_0+C_0+D_0}.
\]
Then $0<\rho_*<1$ and
\begin{equation}\label{eq:product}
 \abs{A_{m,n}^{\omega}}
 \le (C_\omega^-)^{\lambda_*}(C_\omega^+)^{1-\lambda_*}
      \rho_*^{m+n}\norm g_{X_-}\norm f_Y.
\end{equation}
If the interpolated occurrence constants are summable, then
\[
 \sum_\omega\sum_{m,n\ge0}\abs{A_{m,n}^{\omega}}<\infty
\]
uniformly on compact parameter chambers and bounded input sets.
\end{theorem}

\begin{proof}
Raise \eqref{eq:source-side} to the power $\lambda$ and
\eqref{eq:target-side} to the power $1-\lambda$.  Geometric interpolation
gives the time factors
\[
 \alpha^\lambda\Gamma_-^{1-\lambda},
 \qquad
 \Gamma_+^\lambda\beta^{1-\lambda}.
\]
Their logarithms agree at $\lambda_*$ and equal $\log\rho_*$.  This common
number is negative exactly under \eqref{eq:determinant}.  Summation of the two
geometric series and then the occurrence constants proves the result.  The
reverse implication follows because the two strict logarithmic inequalities
have a common solution precisely when their open intervals overlap, which is
again \eqref{eq:determinant}.
\end{proof}

\section{Complete finite differences and higher-source totalization}

For $h\ne0$ let $K_{a,h}=(L_{a+h}-L_a)/h$.

\begin{theorem}[Complete finite-DQ convergence]\label{thm:fdq}
Suppose the finite-DQ response array converges uniformly on each finite
rectangle in $(m,n)$ and that the finite-DQ and limiting arrays share one
summable product envelope independent of small $h$.  Then
\begin{equation}\label{eq:fdq}
 \sum_{m,n\ge0}
 \abs{\pair{Q_{a+h}^nK_{a,h}Q_a^m g-Q_a^nK_aQ_a^m g}{f}}
 \longrightarrow0.
\end{equation}
\end{theorem}

\begin{proof}
Choose a finite rectangle whose common envelope tail is below
$\varepsilon/3$.  Uniform convergence controls the finite head, while the two
tails are bounded by two copies of the envelope tail.  This is convergence in
the complete $\ell^1(\N^2)$ response topology, not merely pointwise
convergence of its coordinates.
\end{proof}

Let $d=(K,\varepsilon,\mathcal P)$ be a depth, intersection, and refinement
cutoff and write
\[
 G_{j,d}(a)=\Tot_{\mathfrak S_d}
 \{\eta^{(j)}_{a,\omega}:\omega\in I_d\}+S_{j,d}^{\rm seam}(a).
\]

\begin{theorem}[Joint-Cauchy higher-source theorem]\label{thm:totalization}
Fix $N\ge1$.  Assume for $1\le j\le N$ that:
\begin{enumerate}[label=(\roman*)]
\item every finite $G_{j,d}$ is continuous in a complete graded source space;
\item $\sup_a\norm{\eta^{(j)}_{a,\omega}}_{\rm CM2}\le b_{j,\omega}$ and
$\sum_\omega b_{j,\omega}<\infty$;
\item intersection and refinement defects are bounded by moduli tending to
zero;
\item common refinements preserve UIDs and cancel opposite orientations;
\item the $j$th letter maps $X_r$ to $X_{r-j}$, while reduced resolvents
preserve every level used by an order-$N$ word.
\end{enumerate}
Then every $G_{j,d}$ is uniformly Cauchy in $a$, and all Kato words of total
order at most $N$ converge to continuous bounded maps between their declared
levels.
\end{theorem}

\begin{proof}
Pass two cutoffs to a common refinement.  Common atoms cancel by UID and
orientation.  The remaining norm is bounded by twice the occurrence tail plus
twice the two defect moduli.  Completeness gives $G_j$.  The graded type of
each letter makes every resolvent word composable, and a finite product of
continuous bounded maps is continuous.
\end{proof}

\section{A symplectic Markov collision map with moving seams}

Let
\[
 I=[-1/40,1/40]
\]
and define
\begin{align*}
 w_1(a)&=1/10+a,& w_2(a)&=1/5+2a,\\
 w_3(a)&=3/10-a,& w_4(a)&=2/5-2a.
\end{align*}
Let $s_0=0$ and $s_i=\sum_{j\le i}w_j$.  On
$I_i(a)=[s_{i-1}(a),s_i(a))$ define
\begin{equation}\label{eq:T}
 T_a(x)={x-s_{i-1}(a)\over w_i(a)}.
\end{equation}
Its two-dimensional collision extension is
\begin{equation}\label{eq:baker}
 B_a(x,y)=\left(T_a(x),\,s_{i-1}(a)+w_i(a)y\right),
 \qquad (x,y)\in I_i(a)\times[0,1].
\end{equation}

\begin{proposition}[Symplectic collision realization]\label{prop:symplectic}
The map $B_a$ is an invertible piecewise affine hyperbolic map of the unit
square modulo its singular seams.  It preserves Lebesgue area, has unstable
multiplier $1/w_i(a)$ and stable multiplier $w_i(a)$ on branch $i$, and its
vertical and horizontal singular seams move nontrivially with $a$.  The
mapping-torus suspension of $B_a$ with any positive bounded roof is a
piecewise smooth finite-flight deterministic collision flow.
\end{proposition}

\begin{proof}
Each rectangle $I_i(a)\times[0,1]$ is mapped bijectively onto
$[0,1]\times I_i(a)$.  The derivative is
\[
 \begin{pmatrix}w_i^{-1}&0\\0&w_i\end{pmatrix},
\]
whose determinant is one.  The four images partition the square, proving
invertibility and area preservation away from the seams.  Hyperbolicity and
the multipliers are immediate.  The mapping torus identifies the top of each
flight cell with its image under $B_a$; the roof bounds give finite flight
times.
\end{proof}

The unstable quotient has Perron--Frobenius operator
\begin{equation}\label{eq:L}
 (L_ah)(y)=\sum_{i=1}^4w_i(a)
 h\bigl(s_{i-1}(a)+w_i(a)y\bigr).
\end{equation}
Lebesgue probability is invariant.  Let $X_r=W^{r,1}(\T)$ and let $X_r^0$
be its zero-mean subspace.  Put
\[
 \rho=\sup_{a\in I,i}w_i(a)=9/20.
\]

\begin{lemma}[Uniform centered contraction]\label{lem:contraction}
For every finite $r\ge1$ there is $C_r$ such that
\[
 \norm{L_a^nh}_{W^{r,1}}
 \le C_r\rho^n\norm h_{W^{r,1}},
 \qquad h\in X_r^0,\quad n\ge0,\quad a\in I.
\]
\end{lemma}

\begin{proof}
For $j\ge1$,
\[
 (L_ah)^{(j)}(y)=
 \sum_iw_i(a)^{j+1}h^{(j)}(s_{i-1}+w_i y).
\]
Changing variables on each branch gives
$\norm{(L_ah)^{(j)}}_1\le\rho^j\norm{h^{(j)}}_1$.  Iterate this estimate.
The zero-mean periodic Poincar\'e inequality controls the $L^1$ norm by the
first derivative.  Summing derivative levels proves the claim.
\end{proof}

Let $v=(1,2,-1,-2)$, $A_i=\sum_{j\le i}v_j$, and
\[
 \psi_i^a(y)=s_{i-1}(a)+w_i(a)y,
 \qquad r_i(y)=A_{i-1}+v_i y.
\]

\begin{lemma}[Exact letters of every order]\label{lem:letters}
For each $k\ge1$,
\begin{equation}\label{eq:letters}
 \partial_a^kL_ah=
 \sum_{i=1}^4\left[
 k v_i r_i^{k-1}h^{(k-1)}(\psi_i^a)
 +w_i(a)r_i^kh^{(k)}(\psi_i^a)
 \right].
\end{equation}
Consequently
\[
 \norm{\partial_a^kL_ah}_{W^{r,1}}
 \le C_{r,k}\norm h_{W^{r+k,1}},
 \qquad \Piop\partial_a^kL_a=0.
\]
\end{lemma}

\begin{proof}
Both $w_i$ and $\psi_i^a$ are affine in $a$.  In the $k$th Leibniz rule,
no derivative or exactly one derivative can hit $w_i$; all remaining
derivatives hit $h(\psi_i^a)$.  This gives \eqref{eq:letters}.  Composition
and product estimates give the norm bound.  Differentiating
$\int L_ah=\int h$ gives centering.
\end{proof}

\begin{theorem}[All-order moving-seam response]\label{thm:allorder}
For every prescribed finite order $N$ the family \eqref{eq:T}--\eqref{eq:L}
has, uniformly in $a\in I$:
\begin{enumerate}[label=(\roman*)]
\item nonzero physical moving-seam currents and a nonzero assembled first
letter;
\item product response bounds
\[
 \abs{\pair{Q_a^n(\partial_a^kL_a)Q_a^m g}{f}}
 \le C_{r,k}\rho^{m+n}
 \norm g_{W^{r+k,1}}\norm f_{(W^{r,1})^*},
 \quad 1\le k\le N+1;
\]
\item finite differences through order $N$ converging in the corresponding
complete $\ell^1(\N^2)$ response topology;
\item $C^N$ Riesz projections, reduced resolvents, pressure eigenvalues, and
all Kato words on the invariant ladder
\[
 W^{r+N,1}\to W^{r+N-1,1}\to\cdots\to W^{r,1};
\]
\item closure under every later insertion on the same ladder, without a
historical-label reset.
\end{enumerate}
\end{theorem}

\begin{proof}
Differentiating the branch indicators before inverse-branch assembly gives
nonzero oriented endpoint currents at $s_1,s_2,s_3$; formula
\eqref{eq:letters} is their complete physical assembly.  Apply
\cref{lem:contraction,lem:letters} on the two sides of every letter to obtain
(ii).  Taylor's formula with the $(N+1)$st letter supplies uniform convergence
on finite time rectangles, and the product bound supplies the common envelope;
\cref{thm:fdq} gives (iii).  The Neumann series
$(I-Q_a)^{-1}=\sum_{n\ge0}Q_a^n$ converges on each centered level.
Repeated differentiation gives only graded products of letters and reduced
resolvents, proving (iv) and (v).
\end{proof}

\begin{theorem}[Genuine nonconjugacy]\label{thm:nonconjugacy}
Distinct parameters in $I$ are not $C^1$ conjugate as collision maps.  In
particular the family is not an exact-coordinate artifact.
\end{theorem}

\begin{proof}
Branches two and three contain regular interior fixed points
\[
 x_i(a)={s_{i-1}(a)\over1-w_i(a)},\qquad i=2,3.
\]
Their unstable multipliers are $1/w_2(a)$ and $1/w_3(a)$.  Over $I$ the two
multiplier ranges are disjoint, so a $C^1$ conjugacy cannot exchange these
fixed points.  A $C^1$ conjugacy preserves the multiplier of each regular
periodic point, while both multipliers vary strictly with $a$.  Therefore two
distinct parameters cannot be conjugate.
\end{proof}

\section{Independently prescribed correlation response}

Let $F_a\in W^{r+N,1}$ be centered and let
$G_a\in(W^{r,1})^*$, both $C^N$ in $a$.  Define
\begin{equation}\label{eq:susceptibility}
 \chi(a;F,G)=
 \sum_{n=0}^\infty\pair{Q_a^nF_a}{G_a}
 =\pair{(I-Q_a)^{-1}F_a}{G_a}.
\end{equation}
No relation of the form $F_a=(I-L_a)H_a$ is assumed.

\begin{theorem}[Arbitrary-source susceptibility]\label{thm:correlation}
For the moving-seam family, \eqref{eq:susceptibility} is $C^N$ in $a$.  Every
$j$th derivative, $j\le N$, is the absolutely convergent sum of all ordered
words made from derivatives of $F_a,G_a$, letters
$\partial_a^kQ_a$, and reduced resolvents, with total parameter order $j$.
The derivative series is normally convergent on $I$.
\end{theorem}

\begin{proof}
The reduced resolvent is a uniformly convergent Neumann series on the centered
ladder.  The resolvent identity gives
\[
 \partial_aR_a=R_a(\partial_aQ_a)R_a,
 \qquad R_a=(I-Q_a)^{-1},
\]
and repeated differentiation gives the stated ordered words.  Each word is
well typed by \cref{lem:letters}; each insertion is surrounded by a product
response envelope from \cref{thm:allorder}.  The finite number of partitions
of $j$ and geometric summability give normal convergence.  Differentiation of
\eqref{eq:susceptibility} is therefore justified term by term.
\end{proof}

This theorem answers a concern that invariant-density differentiation or exact
coboundaries might be the only positive examples: $F$ and $G$ are independently
prescribed, and the response mechanism is the uniform moving-seam dynamics.

\section{Exact predictable innovations}

Let $X_0$ be Lebesgue distributed.  Let
$A_k\in I$ be measurable with respect to
$\F_k=\sigma(I_0,\ldots,I_{k-1})$, where
$I_k=i$ if $X_k\in I_i(A_k)$, and set
\[
 X_{k+1}=T_{A_k}(X_k).
\]

\begin{theorem}[Exact innovation theorem]\label{thm:innovations}
For every $k\ge0$:
\begin{enumerate}[label=(\roman*)]
\item conditionally on $\F_k$, $X_k$ is Lebesgue distributed;
\item $\mathbb P(I_k=i\mid\F_k)=w_i(A_k)$;
\item conditionally on $\F_k\vee\sigma(I_k)$, the image $X_{k+1}$ is
Lebesgue distributed and independent of that sigma-field.
\end{enumerate}
Consequently, for any bounded branch vectors $c_i(a)$,
\[
 M_k=c_{I_k}(A_k)-\sum_iw_i(A_k)c_i(A_k)
\]
is a martingale difference with exact predictable covariance
\[
 \E[M_k\otimes M_k\mid\F_k]
 =\sum_iw_i(A_k)m_i(A_k)\otimes m_i(A_k).
\]
\end{theorem}

\begin{proof}
Induct on $k$.  Conditional on the past, $A_k$ is fixed and $X_k$ is uniform.
The probability of branch $i$ is its length $w_i(A_k)$.  Given that branch,
$X_k$ is uniform on the branch interval, and the affine full-cover map sends
that conditional law to Lebesgue measure independently of the branch label.
This proves the induction and the covariance identity.
\end{proof}

\section{Symbolic desingularization for moving open billiards}

Let $K_1(a),\ldots,K_J(a)\subset\R^d$, $J\ge3$, be a real-analytic family of
pairwise disjoint strictly convex obstacles with analytic boundaries.  Assume
a uniform no-eclipse condition: the convex hull of any two obstacles is
disjoint from every third obstacle.  Let
\[
 \Sigma_A=\{\xi\in\{1,\ldots,J\}^{\Z}:\xi_k\ne\xi_{k+1}\}
\]
with the usual Hölder metric.

\begin{theorem}[Analytic coding on one fixed symbolic space]\label{thm:symbolic}
After shrinking the parameter interval if necessary, there is a uniform
$0<\theta<1$ and a real-analytic map
\[
 a\longmapsto\pi_a\in C^\alpha(\Sigma_A,\mathcal M)
\]
which sends a symbolic itinerary to the corresponding collision state on the
nonwandering set of the open billiard.  The roof
\[
 \tau_a(\xi)=\abs{q_1(a,\xi)-q_0(a,\xi)}
\]
and the unstable Jacobian potential
$\varphi_a=-\log J^uB_a\circ\pi_a$ depend real analytically on $a$ as elements
of $C^\alpha(\Sigma_A)$.  The same is true for every analytic physical
observable pulled back by $\pi_a$.
\end{theorem}

\begin{proof}
For a fixed itinerary, the next collision point is obtained from the current
boundary coordinate and outgoing wavefront by the specular reflection graph
transform.  Strict convexity, positive obstacle separation, and no eclipse
make the inverse graph transform a uniform contraction in a Hölder sequence
norm; the contraction constant is stable on a compact parameter interval.
The analytic parameter-dependent contraction theorem gives the fixed point
$\pi_a(\xi)$ and its analytic dependence on $a$.  Exponential loss of the
distant symbolic coordinates gives the Hölder modulus.  Flight length,
reflection, and unstable-Jacobian formulas are analytic compositions of the
collision point and boundary jets, proving the remaining assertions.
\end{proof}

For a real-analytic Hölder potential $\phi_a$, displacement $\kappa_a$, and
roof $\tau_a$, define on the one-sided shift
\begin{equation}\label{eq:Ruelle}
 (\mathscr L_{a,q,s}f)(\xi)=
 \sum_{\sigma\eta=\xi}
 e^{\phi_a(\eta)+q\cdot\kappa_a(\eta)-s\tau_a(\eta)}f(\eta).
\end{equation}

\begin{theorem}[All-order moving-scatterer response]\label{thm:open-billiard}
For the analytic no-eclipse family, the operator family
\eqref{eq:Ruelle} is analytic on one fixed Hölder space.  In a neighbourhood
of $(q,s)=(0,0)$ its leading eigenvalue, Riesz projection, equilibrium state,
reduced resolvent, and pressure are analytic in $(a,q,s)$.  For arbitrary
analytic Hölder observables $F_a,G_a$ the correlation susceptibility
\[
 \sum_{n\ge0}
 \left(\mu_a(F_aG_a\circ\sigma^n)-\mu_aF_a\,\mu_aG_a\right)
\]
is analytic in $a$, and every finite derivative is given by a normally
convergent Ruelle--resolvent word.  These claims apply to the corresponding
physical collision observables through the coding map $\pi_a$.
\end{theorem}

\begin{proof}
By \cref{thm:symbolic}, every exponent in \eqref{eq:Ruelle} is analytic as a
Hölder function.  The exponential Nemytskii map and the finite preimage sum
make $\mathscr L_{a,q,s}$ an analytic operator family.  Ruelle's theorem gives
a simple isolated eigenvalue and a spectral gap; compactness of the parameter
interval makes the contour uniform.  Analytic Riesz calculus yields the
spectral data.  The centered correlations decay exponentially uniformly.
Their parameter derivatives are finite sums of resolvent words with the same
uniform exponential envelope, hence the susceptibility and all its finite
derivatives converge normally.
\end{proof}

\begin{corollary}[A genuinely nonconjugate specular family]\label{cor:specular}
Choose an asymmetric no-eclipse configuration with a unique shortest regular
two-periodic orbit between $K_1$ and $K_2$.  Move $K_1$ analytically in the
normal direction while keeping the remaining obstacles fixed.  For a small
parameter interval the no-eclipse hypotheses persist, the two-periodic flight
length has nonzero derivative, and the family is not time-preservingly
conjugate.  Nevertheless all conclusions of \cref{thm:open-billiard} hold.
Moreover one may choose a locally constant symbolic observable whose periodic
sums are not constant, so the positive correlation theorem is not an
exact-coboundary benchmark.
\end{corollary}

\begin{proof}
The isolated flight persists by the implicit-function theorem, and its length
derivative is twice the normal velocity projected onto the flight direction,
which is nonzero by construction.  An isolated periodic-flow length is
preserved by a time-preserving conjugacy, proving nonconjugacy.  The no-eclipse
condition and strict curvature are open.  Finally, a cylinder function taking
different sums on two periodic symbolic orbits cannot be cohomologous to a
constant by the periodic-orbit criterion.
\end{proof}

\section{Conclusion}

The two positive models close complementary parts of the moving-singularity
problem.  The symplectic Markov collision map retains physical moving seams and
proves all-order product response, arbitrary-source susceptibility, repeated
insertion, and exact innovations by explicit estimates.  The open dispersing
billiard uses symbolic desingularization to put a genuinely moving specular
system on one fixed Hölder space and thereby proves all-order response for
arbitrary Hölder observables.  In both cases nonconjugacy is certified by a
periodic invariant, and neither proof relies on differentiated invariance or
an exact-coboundary ansatz.

\section*{Disclosure and review status}
The author is responsible for all proofs, references, and computations.
Finite-dimensional formula checks and repository verifiers do not certify the
analytical arguments.  This revision is prepared for a new independent formal
review.

\bibliographystyle{plain}
\bibliography{references-round4}

\end{document}
